\chapter{Markov Chains and Mixing}

\begin{center}\itshape
Move the top card of a deck to a random position, over and over.\\
Once the original bottom card has been moved, the deck is uniform.
\end{center}

A random walk on a graph steps to a neighbour. A Markov chain on a finite set may favour some edges. The question is whether the distribution forgets its start.

\section{Transition matrices}

Let $\mathcal X$ be a finite set. A transition matrix satisfies $P(x,y)\ge0$ and $\sum_y P(x,y)=1$. A chain with initial law $\mu$ is built by drawing each next state from the current row, using fresh blocks of bits. The law of $X_n$ is $\mu P^n$.

\begin{examplex}[three states]
On $\{1,2,3\}$, let
\[
 P=\begin{pmatrix}1/2&1/2&0\\1/3&1/3&1/3\\0&1/2&1/2\end{pmatrix}.
\]
The vector $\pi=(2/7,3/7,2/7)$ satisfies $\pi P=\pi$. The first coordinate of $\pi P$ is $(2/7)(1/2)+(3/7)(1/3)=2/7$, and the other two match by the same arithmetic.
\end{examplex}

\section{Irreducibility and stationary laws}

The chain is \emph{irreducible} if for every pair $x,y$ some $P^n(x,y)$ is positive. The \emph{period} of a state is the greatest common divisor of the return times. In an irreducible chain the period is the same for every state. The chain is \emph{aperiodic} when the period is $1$.

\begin{theorem}[existence and uniqueness]\label{thm:stationary}
Every finite transition matrix has a stationary distribution. An irreducible chain has exactly one.
\end{theorem}
\begin{proof}
From any start, the averages $\nu_N=N^{-1}\sum_{n=0}^{N-1}\mu_n$ lie in the simplex. Bolzano--Weierstrass in each coordinate, Volume~1, Chapter~3, and a diagonal extraction produce a limit $\pi$. Then $\nu_N P-\nu_N=N^{-1}(\mu_N-\mu_0)\to0$, so $\pi P=\pi$.

If $\pi$ and $\pi'$ are stationary and the chain is irreducible, a positive maximum of $\pi-\pi'$ at $x$ propagates along every edge with $P(x,y)>0$, because $(\pi-\pi')(x)$ equals the $P(x,\cdot)$-average of $\pi-\pi'$. Irreducibility forces $\pi-\pi'$ constant, hence zero.
\end{proof}

\section{Convergence}

Total variation is $\|\mu-\nu\|=\frac12\sum_x|\mu(x)-\nu(x)|=\max_A|\mu(A)-\nu(A)|$.

\begin{theorem}[convergence]\label{thm:ergodic}
If $P$ is irreducible and aperiodic on a finite set, with stationary law $\pi$, then $\|P^n(x,\cdot)-\pi\|\to0$ for every $x$.
\end{theorem}
\begin{proof}
Some $n_0$ has $P^{n_0}(u,v)>0$ for all $u,v$. Let $\delta=\min P^{n_0}>0$. Couple two chains, one from $x$ and one from $\pi$, in blocks of length $n_0$. From current states $u$ and $v$, with probability $\delta$ send both to a common draw, and with probability $1-\delta$ move them independently under the residual kernels $(P^{n_0}(u,\cdot)-\delta)/(1-\delta)$. They meet with probability at least $\delta$ on each block, and stay together afterwards. The probability they differ at time $n$ is at most $(1-\delta)^{\lfloor n/n_0\rfloor}$, which tends to $0$, and total variation is at most that probability.
\end{proof}

For random transpositions the sharp window sits near $\frac12 n\ln n$ and is proved in Volume~11. The bound above is only exponential convergence.

\section{Reversibility}

The chain is reversible for $\pi$ when $\pi(x)P(x,y)=\pi(y)P(y,x)$. Any such $\pi$ is stationary. The walk on an undirected graph that steps to a uniform neighbour is reversible for $\pi$ proportional to degree.

\section{Exercises}

\begin{exercise}[$\star$]\label{ex:lazy}
Show that $Q=(P+I)/2$ is aperiodic, has the same stationary distributions as $P$, and is irreducible whenever $P$ is.
\end{exercise}

\begin{exercise}[$\star$]\label{ex:twostate}
On $\{0,1\}$ let $P(0,1)=p$ and $P(1,0)=q$ with $p,q\in(0,1]$. Find the stationary distribution and $\|\mathcal L(X_n)-\pi\|$ starting from $0$.
\end{exercise}

\begin{exercise}[$\star\star$]\label{ex:ehrenfest}
The Ehrenfest chain on $\{0,\ldots,d\}$ moves from $k$ to $k+1$ with probability $(d-k)/d$ and to $k-1$ with probability $k/d$. Show that $\pi(k)=\binom{d}{k}2^{-d}$ is stationary.
\end{exercise}

\begin{exercise}[$\star\star$]\label{ex:pagerank}
From a directed graph with minimum out-degree $1$, follow a uniform out-edge with probability $\alpha\in(0,1)$ and otherwise jump to a uniform vertex. Show that the chain has a unique stationary distribution.
\end{exercise}

\begin{exercise}[$\star\star$]\label{ex:metropolis}
Let $\pi>0$ and let $Q$ be symmetric and irreducible. Set $P(x,y)=Q(x,y)\min(1,\pi(y)/\pi(x))$ for $y\ne x$, with the rest on the diagonal. Show that $P$ is reversible for $\pi$.
\end{exercise}

\begin{exercise}[$\star\star$]\label{ex:toptorandom}
Move the top card of an $n$-card deck to a uniform random position. Let $\tau$ be the first time the original bottom card has been inserted. Show that the deck at time $\tau$ is uniform and independent of $\tau$, that $\tau$ is a coupon-collector time, and that the total variation distance to uniform at time $n\ln n+cn$ is at most $e^{-c}$.
\end{exercise}

\begin{exercise}[$\star$]\label{ex:period}
Show that the cycle $1\to2\to\cdots\to n\to1$ is irreducible of period $n$, with unique stationary law the uniform law, and that $P^n$ does not converge to that law if $n>1$.
\end{exercise}

\begin{exercise}[$\star\star$]\label{ex:risingmean}
The number of rising sequences in a permutation is one plus the number of descents, positions $i$ where the card in place $i$ exceeds the card in place $i+1$. Show that a uniform permutation of $n\ge2$ cards has mean $(n+1)/2$ and variance $(n+1)/12$ rising sequences. For $n=52$, deduce that more than $90\%$ of permutations have between $20$ and $33$ rising sequences.
\end{exercise}

\section{Hints}

\begin{hint}{ex:lazy}
A positive diagonal kills the period.
\end{hint}
\begin{hint}{ex:twostate}
$\pi=(q,p)/(p+q)$, and the error multiplies by $1-p-q$ at each step.
\end{hint}
\begin{hint}{ex:ehrenfest}
Check $\pi(k)P(k,k+1)=\pi(k+1)P(k+1,k)$.
\end{hint}
\begin{hint}{ex:pagerank}
Every entry of $P$ is at least $(1-\alpha)/|\mathcal X|$.
\end{hint}
\begin{hint}{ex:metropolis}
$\pi(x)P(x,y)=Q(x,y)\min(\pi(x),\pi(y))$.
\end{hint}
\begin{hint}{ex:toptorandom}
While $k$ cards lie below the original bottom card, the next insertion falls below it with probability $(k+1)/n$. The tail of the coupon time is at most $e^{-c}$.
\end{hint}
\begin{hint}{ex:period}
$P^n=I$.
\end{hint}
\begin{hint}{ex:risingmean}
Non-adjacent descents are independent. Adjacent ones occur together with probability $1/6$.
\end{hint}

\section{Solutions}

\begin{solution}{ex:lazy}
\writeup $Q(x,x)\ge1/2$, so the period is $1$. Edges of $P$ remain edges of $Q$. Stationarity is equivalent for $P$ and $Q$.
\end{solution}

\begin{solution}{ex:twostate}
\writeup Balance gives $\pi_0=q/(p+q)$ and $\pi_1=p/(p+q)$. The difference $\delta_n=P(X_n=0)-\pi_0$ satisfies $\delta_{n+1}=(1-p-q)\delta_n$. From state $0$, $\delta_0=p/(p+q)$, so the total variation distance is $|\delta_n|$.
\end{solution}

\begin{solution}{ex:ehrenfest}
\writeup $\pi(k)(d-k)/d=\binom{d}{k}2^{-d}(d-k)/d=\binom{d}{k+1}2^{-d}(k+1)/d=\pi(k+1)P(k+1,k)$.
\end{solution}

\begin{solution}{ex:pagerank}
\writeup Every entry is positive, so the chain is irreducible and aperiodic. Theorems~\ref{thm:stationary} and~\ref{thm:ergodic} apply.
\end{solution}

\begin{solution}{ex:metropolis}
\writeup For $y\ne x$, both $\pi(x)P(x,y)$ and $\pi(y)P(y,x)$ equal $Q(x,y)\min(\pi(x),\pi(y))$, because $Q$ is symmetric.
\end{solution}

\begin{solution}{ex:toptorandom}
\writeup Let $B$ be the original bottom card. Cards below $B$ stay in their relative order until $B$ itself is moved. The number below $B$ increases as in the coupon collector: with $k$ cards already below, the top card is inserted below $B$ with probability $(k+1)/n$. The time $\tau$ until $B$ has been inserted is the collector's time $T$ of Theorem~\ref{thm:coupon}. At that time the order is a uniform permutation, independent of $\tau$, because each insertion that built the order was uniform. The bound $P(T>n\ln n+cn)\le n(1-1/n)^{n\ln n+cn}\le e^{-c}$ is the estimate in the proof of that theorem's tail, using $1-1/n\le e^{-1/n}$. Total variation at a fixed time $t$ is at most $P(\tau>t)$.
\end{solution}

\begin{solution}{ex:period}
\writeup The chain reaches every state, returns only at multiples of $n$, and preserves the uniform law. For $n>1$, $P^n=I$, so a point mass returns to itself.
\end{solution}

\begin{solution}{ex:risingmean}
\writeup Let $I_j$ indicate a descent at $j=1,\ldots,n-1$, and let $D=\sum I_j$, so the number of rising sequences is $R=D+1$. Each pair of adjacent positions is in either order with probability $1/2$, hence $\mathbb E[D]=(n-1)/2$ and $\mathbb E[R]=(n+1)/2$. If $|j-k|\ge2$ the pairs are disjoint and independent, so the covariance vanishes. If the descents are adjacent, three cards have $6$ orders and only one is a double descent, so the joint probability is $1/6$ and the covariance is $1/6-1/4=-1/12$. There are $2(n-2)$ ordered adjacent pairs in the variance sum, and
\[
 \operatorname{Var}D=\frac{n-1}{4}-\frac{n-2}{6}=\frac{n+1}{12}.
\]
For $n=52$, $\mathbb E[R]=26.5$ and $\operatorname{Var}R=53/12$. Chebyshev's inequality at distance $7.5$ gives
\[
 P(|R-26.5|\ge 7.5)\le\frac{53/12}{56.25}=\frac{53}{675}<0.079,
\]
so more than $92\%$ of permutations have $R\in\{20,\ldots,33\}$.
\end{solution}

\section*{Sources and further reading}

\begin{enumerate}
\item A.~A.~Markov, Extension of the limit theorems of probability theory to a sum of variables connected in a chain, 1906.
\item W.~Doeblin, Expos\'e de la th\'eorie des cha\^ines simples constantes de Markoff \`a un nombre fini d'\'etats, \emph{Rev.\ Math.\ Union Interbalkan.} \textbf{2} (1938), 77--105.
\end{enumerate}

